\begin{figure*}[t]
  \centering
  \resizebox{0.97\textwidth}{!}{%
  \begin{tikzpicture}[
    font=\sffamily\scriptsize,
    >=Latex,
    panel/.style={draw=black!45, rounded corners=2pt, fill=black!1},
    good/.style={circle, draw=blue!70!black, fill=blue!18, inner sep=1.5pt},
    bad/.style={circle, draw=red!70!black, fill=red!18, inner sep=1.5pt},
    proto/.style={circle, draw=orange!80!black, fill=orange!22, inner sep=1.6pt},
    mapped/.style={rectangle, draw=green!50!black, fill=green!20, inner sep=1.5pt},
    flow/.style={-{Latex[length=2.2mm]}, line width=0.75pt},
    thinflow/.style={-{Latex[length=1.7mm]}, line width=0.55pt}
  ]

  % Panel A: single direction
  \begin{scope}[xshift=0cm]
    \draw[panel] (0,0) rectangle (6.6,4.65);
    \node[font=\sffamily\bfseries\small] at (3.3,4.34) {(a) 单方向投影};
    \node[good] at (1.15,1.20) {};
    \node[good] at (1.55,1.62) {};
    \node[good] at (1.85,0.93) {};
    \node[good] at (2.15,1.42) {};
    \node[bad] at (4.25,2.88) {};
    \node[bad] at (4.65,3.24) {};
    \node[bad] at (4.92,2.55) {};
    \node[bad] at (5.25,3.02) {};
    \coordinate (mug) at (1.7,1.3);
    \coordinate (mub) at (4.75,2.92);
    \fill[blue!70!black] (mug) circle (2pt) node[below=2pt] {$\mu_g$};
    \fill[red!70!black] (mub) circle (2pt) node[above=2pt] {$\mu_b$};
    \draw[flow, purple!75!black] (mug) -- node[above, sloped] {$r=\mu_b-\mu_g$} (mub);
    \draw[dashed, black!55] (0.7,2.75) -- (5.9,0.45) node[below left] {$r^\perp$};
    \draw[thinflow, red!70!black] (4.25,2.88) -- (3.37,1.85);
    \draw[thinflow, red!70!black] (4.92,2.55) -- (3.70,1.43);
    \node[align=center] at (3.3,0.32) {删除一个全局分量：$x\mapsto\Pi_r(x)$};
  \end{scope}

  % Panel B: SOM directions
  \begin{scope}[xshift=7.05cm]
    \draw[panel] (0,0) rectangle (7.35,4.65);
    \node[font=\sffamily\bfseries\small] at (3.675,4.34) {(b) SOM 多方向};
    \draw[step=0.72cm, black!35] (0.85,0.75) grid (3.73,3.63);
    \foreach \i in {0,...,3}{
      \foreach \j in {0,...,3}{
        \node[proto] at ({0.85+0.72*\i},{0.75+0.72*\j}) {};
      }
    }
    \fill[blue!70!black] (5.80,1.10) circle (2.2pt) node[below=2pt] {$\nu$};
    \draw[thinflow, orange!80!black] (5.80,1.10) -- node[below, sloped] {$r_1$} (1.57,1.47);
    \draw[thinflow, orange!80!black] (5.80,1.10) -- node[above, sloped] {$r_2$} (3.01,2.91);
    \draw[thinflow, orange!80!black] (5.80,1.10) -- node[above, sloped] {$r_3$} (1.57,3.63);
    \node[align=center] at (5.78,2.55) {$r_i=w_i-\nu$\\16 个候选};
    \node[align=center] at (3.675,0.30) {BO/TPE 选择有序的 $k=2,\ldots,7$ 个投影};
  \end{scope}

  % Panel C: local matrix mapping
  \begin{scope}[xshift=14.85cm]
    \draw[panel] (0,0) rectangle (8.1,4.65);
    \node[font=\sffamily\bfseries\small] at (4.05,4.34) {(c) 上游局部矩阵映射};
    \node[bad] (xb1) at (0.75,2.75) {};
    \node[bad] (xb2) at (0.75,2.15) {};
    \node[bad] (xb3) at (0.75,1.55) {};
    \node[align=center] at (0.75,0.93) {模块输入\\$X_b$};
    \node[draw, rounded corners=2pt, fill=purple!12, minimum width=1.05cm, minimum height=1.05cm] (W) at (2.25,2.15) {$W$};
    \draw[flow] (1.02,2.15) -- (W);
    \node[mapped] (z1) at (4.05,2.80) {};
    \node[mapped] (z2) at (4.38,2.30) {};
    \node[mapped] (z3) at (3.82,1.82) {};
    \draw[flow, purple!75!black] (W) -- node[above] {$X_bW^\top$} (3.62,2.30);
    \node[good] at (5.15,2.92) {};
    \node[good] at (5.52,2.48) {};
    \node[good] at (5.12,1.95) {};
    \node[good] at (5.72,2.02) {};
    \node[align=center, blue!70!black] at (5.45,3.52) {良性输出邻域 $Y_g$};
    \node[bad] at (7.15,1.10) {};
    \node[bad] at (7.50,1.48) {};
    \node[bad] at (7.20,1.86) {};
    \node[align=center, red!70!black] at (7.18,0.48) {原目标输出 $Y_b$};
    \draw[thinflow, green!50!black] (4.38,2.30) -- node[above, sloped] {拉近} (5.12,2.48);
    \draw[thinflow, red!70!black, dashed] (4.05,1.82) -- node[below, sloped] {推离} (6.98,1.30);
    \node[align=center] at (3.95,0.36) {直接优化映射；不显式构造 $r_i$};
  \end{scope}

  \node[font=\sffamily\footnotesize, align=center] at (11.45,-0.52)
    {三种参数化的结构并列；面板顺序不表示理论包含关系，也不表示效果递增。};
  \end{tikzpicture}%
  }
  \caption{从显式方向到直接映射的几何对照。蓝色表示良性参照，红色表示目标行为样本，橙色表示 SOM 原型，绿色方块表示矩阵映射后的代理输出。SOM-MD 仍以方向和投影为显式变量；面板 (c) 示意上游矩阵代理；本地 point-v1 使用缓存输出加低秩增量与软邻域。该图只比较结构，不表示实测轨迹、机制因果性或跨协议性能。}
  \label{fig:method-geometry}
\end{figure*}
